%% =========================================================================
\section{The Danish Building and Housing Register (BBR)}\label{sec:bbr}

\draftnote{This section is descriptive and drafted in full: it renders
the register facts we verified against the official BBR code lists, the
grunddata object model, and the BBR case-processing instructions
(docs/indicator-decision.md and docs/indicators.md), plus the published
BUILD reports. Please check the wording against the sources; nothing
here should be new analysis.}

\subsection{Purpose, structure, and distribution}\label{sec:bbr-structure}

The Danish Building and Housing Register (Bygnings- og Boligregistret,
BBR) is the national administrative register of buildings and dwellings,
established by the Building and Housing Registration Act of 1976
(BBR-loven)~\cite{bbrloven}. It is maintained by the municipalities, and
%% NOTE: the register became operational/was populated in 1977; add that
%% detail only with a verified source if the group wants it.
property owners are obliged to report changes affecting their buildings.
The register is distributed through the national data portal, the Danish
Data Distributor (Datafordeler), as part of the Danish common
public-sector data program~\cite{datafordeler_grunddatamodel}.

The central object for this study is the building (\texttt{Bygning}).
Records are versioned and bitemporal~\cite{bbr_bitemporalitet}: each
version carries registration timestamps (\emph{registreringstid}, when
the version was recorded in the register) and effect timestamps
(\emph{virkningstid}, \texttt{virkningFra}/\texttt{virkningTil}, when
the described state applied). A building's register history is a
sequence of such versions. Datafordeler distributes both a
current-state view of the register and extracts that include the full
version history, i.e., all versions and all lifecycle statuses.

Two properties of this distribution matter for demolition analysis.
First, a completed demolition removes the building from the
current-state view, so demolished buildings are absent from a current
snapshot; analyses must use the history-including extract, in which the
building is retained with a lifecycle status of
``historical''~\cite{datafordeler_grunddatamodel}. Second, the
register's history on Datafordeler begins at the conversion of BBR
(version 1.8) to the common data model: all data converted at that
point carry the registration timestamp 2~June~2017, and the official
documentation warns that queries with registration or effect times
before that date may return missing or incomplete
data~\cite{bbr_bitemporalitet}.

\subsection{How a demolition is recorded}\label{sec:bbr-chain}

Alongside buildings, BBR registers building cases (byggesager): the
administrative cases under which municipalities process construction,
alteration, and demolition matters. A case is linked to the buildings
it concerns and carries a case type, case dates, and a case
status~\cite{bbr_instruks, datafordeler_grunddatamodel}. BBR's
case-processing instructions define a demolition workflow on this
basis~\cite{bbr_instruks}. A demolition is handled as a building case
of the demolition type: historically case type (felt 290) 3,
``Nedrivning (hel eller delvis)'', which the current distributed code
list splits into \texttt{sagstype} 31 (partial) and 32
(total)~\cite{bbr_kodelister}. The case carries a notification date
(felt 294, ``Anmeldelse af nedrivning''). When the demolition is
reported physically completed, the completion date (felt 295,
``Gennemf{\o}rt nedrivning'') is recorded; reporting felt 295 removes
the building from the current stock, and its lifecycle status becomes
10, ``Historisk''~\cite{bbr_instruks}.

The completion fields themselves, however, are not part of the
distributed data: felt 294 and felt 295 exist in the operational
register but are not exposed as attributes of the distributed building
or case objects~\cite{datafordeler_grunddatamodel}. What the public
distribution carries instead are two traces of the workflow: the
lifecycle-status change of the building, and the linked demolition case.
Published register-based studies have accordingly worked with one or
both of these traces, or with extracts compiled by parties with access
to the operational register (Section~\ref{sec:indicators}).

\subsection{Observable demolition-related attributes}\label{sec:bbr-fields}

Three distributed attribute families carry demolition-related
information. Code semantics below are taken from the official BBR code
lists~\cite{bbr_kodelister} and the grunddata object-type
catalog~\cite{datafordeler_grunddatamodel}.

\paragraph{Lifecycle status (\texttt{Bygning.status}).}
Table~\ref{tab:status-codes} lists the relevant codes. Status 10
(``Historisk'') is the state a completed demolition leads to. The code
list also defines a separate terminal state for erroneous registrations,
status 11 (``Fejlregistreret'').

\begin{table}[htbp]
\centering
\caption{Lifecycle (Livscyklus) codes of the distributed building object
relevant to demolition analysis~\cite{bbr_kodelister}.}
\label{tab:status-codes}
\begin{tabular}{@{}lll@{}}
\toprule
Code & Label (Danish) & Meaning \\
\midrule
6  & Opf{\o}rt            & Completed, in active stock \\
7  & G{\ae}ldende         & Current/valid, in active stock \\
10 & Historisk            & Historical: out of the active stock \\
11 & Fejlregistreret      & Erroneous registration (separate terminal state) \\
13 & Delvis Afsluttet     & Partially completed \\
\bottomrule
\end{tabular}
\end{table}

\paragraph{Business process (\texttt{Bygning.forretningsproces}).}
Each record version carries a code stating the business event that
caused the update; Table~\ref{tab:bp-codes} lists the codes relevant
here. Code 3 (``Opdateret grundet nedrivning'') marks an update due to
demolition. The list also names several non-demolition causes of record
updates, including error correction (4), splits and merges (6, 7),
synchronization with other base registers (12--15), and cadastral
restructuring (17).

\begin{table}[htbp]
\centering
\caption{Business-process (ForretningsProcess) codes relevant to
demolition analysis~\cite{bbr_kodelister}.}
\label{tab:bp-codes}
\begin{tabular}{@{}lp{6.2cm}l@{}}
\toprule
Code & Label (Danish, abbreviated) & Event \\
\midrule
1      & Oprettet grundet nybyggeri & Creation (new construction) \\
2      & Opdateret grundet til/ombygning & Renovation/extension \\
3      & Opdateret grundet nedrivning & Update due to demolition \\
4      & Fejlrettelse af faktiske fejlregistreringer & Error correction \\
5      & Faktisk udf{\o}rt {\ae}ndring uden byggesagsbehandling & Change without building case \\
6, 7   & Opdeling / Sammenl{\ae}gning af enheder & Split / merge \\
12--15 & Opdateret grundet {\ae}ndring i grunddataregister & Base-register synchronization \\
17     & Matrikul{\ae}r {\ae}ndring & Cadastral restructuring \\
\bottomrule
\end{tabular}
\end{table}

\paragraph{Building cases.}
Distributed case objects link buildings to building cases. Two codings
identify demolition cases: the current case-type list
(\texttt{Sagsniveau.sagstype} 31 partial / 32 total demolition) and the
older regulatory case code (\texttt{BBRSag.sag012Byggesagskode}), where
code 6 (``Tilladelsessag Nedrivning'') is the demolition permit case and
legacy codes 2 and 5 are marked as being phased out; legacy code 2
covers demolition notifications together with minor-structure
notifications (garages, carports, sheds)~\cite{bbr_kodelister}.

\subsection{Floor-area attributes}\label{sec:bbr-areas}

The building object carries several area
attributes~\cite{datafordeler_grunddatamodel}:

\begin{itemize}
  \item \texttt{byg038SamletBygningsareal} --- total building area (sum
        of floor areas across storeys);
  \item \texttt{byg041BebyggetAreal} --- built-up area (ground-floor
        footprint);
  \item \texttt{byg039BygningensSamledeBoligAreal} and
        \texttt{byg040BygningensSamledeErhvervsAreal} --- the
        residential and commercial floor-area components. Their sum is
        the area definition used in BUILD's national demolition
        analysis~\cite{jensen2025nedrivning};
  \item \texttt{byg054AntalEtager} --- the number of storeys.
\end{itemize}

The completeness of each attribute in the current extract is reported in
Section~\ref{sec:missingness}. \draftnote{Per our data rules, the
missingness figures we know from earlier work (byg038 $\approx$44\,\%
missing, byg041 $\approx$0.01\,\%) are not stated here --- they are
recomputed on the fresh extract and enter the paper only from
\texttt{results/}.}

\subsection{Data-quality context}\label{sec:bbr-quality}

Register information is reported by property owners and processed by
municipalities \refneeded{Rigsrevisionen report on BBR data quality; DR
(2023) survey --- to be verified and added, or dropped}. The register
has been through system migrations, most recently the 2017 conversion to
the common data model (Section~\ref{sec:bbr-structure}). BUILD's
analysis of retired buildings notes that retirement from the register is
presumed, but not certain, to mean demolition; that retired areas are
probably overestimated; and that buildings registered as retired but
constructed after 1998 largely could not be confirmed as demolished in
manual spot checks, which led BUILD to restrict their analysis to
buildings constructed before 1999~\cite{jensen2025nedrivning}.
